\documentclass{amsart}
\usepackage{amsmath,amssymb,amsthm,hyperref}
\newtheorem{theorem}{Theorem}
\newtheorem{lemma}{Lemma}
\newtheorem{proposition}{Proposition}
\newtheorem{corollary}{Corollary}
\theoremstyle{definition}
\newtheorem{definition}{Definition}
\newtheorem{remark}{Remark}
\newtheorem{example}{Example}
\begin{document}

\title{Connected signed graphs with two distinct eigenvalues}
\maketitle

\section{Setup, conventions and statement of the result}

Throughout, $(G,\sigma)$ is a signed graph with $G=(V,E)$ a finite simple
graph on $n=|V|\ge 2$ vertices and $\sigma\colon E\to\{+1,-1\}$ a signature.
We write $A=A^\sigma$ for the symmetric $n\times n$ adjacency matrix:
$A_{uv}=\sigma(uv)$ if $uv\in E$ and $A_{uv}=0$ otherwise; in particular
$A_{uu}=0$. As $A$ is real symmetric it has $n$ real eigenvalues counted with
multiplicity. We assume $G$ is connected.

\paragraph{Switching.}
For $D=\operatorname{diag}(\varepsilon_1,\dots,\varepsilon_n)$ with
$\varepsilon_i\in\{+1,-1\}$, the matrix $D A D$ is again the adjacency matrix
of a signed graph on the same underlying graph $G$, with signature
$\sigma'(uv)=\varepsilon_u\varepsilon_v\sigma(uv)$. Two signed graphs related
this way are \emph{switching equivalent}. Since $D=D^{-1}=D^\top$, $DAD$ is
similar to $A$, so switching equivalent signed graphs have the same spectrum;
switching also preserves connectedness. We classify up to switching
equivalence and isomorphism, as required.

\paragraph{Notation.} $I$ is the identity matrix, $J$ the all-ones matrix,
$\mathbf 1$ the all-ones vector. We repeatedly use
\begin{equation}\label{eq:walk}
(A^2)_{ij}=\sum_{h=1}^n A_{ih}A_{hj}\quad(i\ne j),\qquad
(A^2)_{ii}=\sum_{h=1}^n A_{ih}^2=\deg_G(i),
\end{equation}
the last equality because $A_{ih}^2=1$ exactly when $ih\in E$ and $0$ otherwise.

\begin{definition}[Symmetric weighing matrix]\label{def:weigh}
A \emph{symmetric weighing matrix of weight $k$ and order $n$ with zero
diagonal} is a symmetric $n\times n$ matrix $W$ with entries in
$\{-1,0,+1\}$, zero diagonal, and $W^2=kI$. Equivalently it is the adjacency
matrix of a signed graph in which every vertex has degree $k$ and $A^2=kI$. We
write $W(n,k)$ for the class of such matrices.
\end{definition}

\begin{definition}[Regular two-graph; Seidel matrix]\label{def:twograph}
A \emph{Seidel matrix} is a symmetric $\{0,\pm1\}$ matrix with zero diagonal
in which \emph{every} off-diagonal entry is $\pm1$; equivalently it is the
adjacency matrix $A^\sigma$ of a signed graph whose underlying graph is the
\emph{complete} graph $K_n$. A Seidel matrix $S$ is called a \emph{regular
two-graph} (we use this spectral definition; it is one of the standard
equivalent definitions) if $S$ has exactly two distinct eigenvalues. Two
regular two-graphs are equivalent precisely when their Seidel matrices are
switching equivalent.
\end{definition}

\begin{theorem}\label{thm:main}
Let $(G,\sigma)$ be a connected signed graph on $n\ge 2$ vertices with
adjacency matrix $A$. Then $A$ has exactly two distinct eigenvalues $a>b$ if
and only if there is an integer $k\ge1$ such that $G$ is $k$-regular and
\begin{equation}\label{eq:quad}
A^2=(a+b)\,A+k\,I,\qquad k=-ab .
\end{equation}
Writing $s=a+b$, exactly one of the following holds, the partition being
governed by the sign of $s$.

\smallskip
\textbf{(I) Balanced case $s=0$.}
Then $a=\sqrt k$, $b=-\sqrt k$, both eigenvalues have multiplicity $n/2$ (so
$n$ is even), and \eqref{eq:quad} reads $A^2=kI$. Thus $A\in W(n,k)$ is a
connected symmetric weighing matrix of weight $k$ with zero diagonal;
conversely every connected $A\in W(n,k)$ is the adjacency matrix of a signed
graph with the two eigenvalues $\pm\sqrt k$. This is an exact
correspondence:
\[
\Bigl\{\text{connected two-eigenvalue signed graphs with }s=0\Bigr\}
\;=\;\bigcup_{k\ge1}\Bigl\{\text{connected }W(n,k)\Bigr\}.
\]

\smallskip
\textbf{(II) Integral case $s\ne 0$.}
Then $a$ and $b$ are integers with $a>0>b$, and, writing $p:=a\ge1$,
$q:=-b\ge1$ (so $p\ne q$, $k=pq$), the underlying graph is the complete graph
$K_n$; equivalently $A$ is a Seidel matrix, i.e.\ a regular two-graph with the
two eigenvalues $p$ and $-q$. Conversely every regular two-graph with $s\ne0$
is such a signed graph. Within this class:
\begin{itemize}
\item[(IIa)] If $a$ has multiplicity $1$ then, up to switching, $A=J-I$;
$(G,\sigma)$ is switching equivalent to the all-positive $K_n$, with spectrum
$\{\,n-1,\ \underbrace{-1,\dots,-1}_{n-1}\,\}$, i.e.\ $a=n-1$, $b=-1$.
\item[(IIb)] If $b$ has multiplicity $1$ then, up to switching, $A=-(J-I)$;
$(G,\sigma)$ is switching equivalent to the all-negative $K_n$, with spectrum
$\{\,\underbrace{1,\dots,1}_{n-1},\ -(n-1)\,\}$, i.e.\ $a=1$, $b=-(n-1)$.
\item[(IIc)] If both multiplicities are $\ge2$ then $(G,\sigma)$ is a regular
two-graph that is \emph{not} switching equivalent to $\pm(J-I)$. This case is
non-empty: the smallest instance is on $n=16$ vertices with spectrum
$\{3^{(10)},(-5)^{(6)}\}$ (Example~\ref{ex:rook}). The complete classification
of this case coincides with the classification of regular two-graphs with two
multiplicities $\ge2$, which is a deep and (in full generality) open problem of
algebraic combinatorics.
\end{itemize}
\end{theorem}

The structure of the paper is as follows. Section~\ref{sec:quad} proves the
equivalence with \eqref{eq:quad} and regularity; Section~\ref{sec:dichotomy}
establishes the dichotomy $s=0$ versus $s\ne0$ and the integrality in
case~(II); Section~\ref{sec:complete} proves that $s\ne0$ forces $G=K_n$, i.e.\
identifies case~(II) with regular two-graphs, and settles the simple-eigenvalue
subcases (IIa),(IIb); Section~\ref{sec:examples} records the converses, the
non-emptiness of all surviving families, and---crucially---the explicit
counterexample showing that case~(IIc) really occurs, so that case~(II) is
\emph{not} reducible to the switched complete graphs.

\section{The defining quadratic equation and regularity}\label{sec:quad}

\begin{lemma}\label{lem:quad}
For two distinct reals $a>b$, the matrix $A$ has eigenvalues lying in
$\{a,b\}$ and uses both values if and only if $(A-aI)(A-bI)=0$, equivalently
\begin{equation}\label{eq:quad2}
A^2=(a+b)A-ab\,I .
\end{equation}
In particular $A$ has exactly two distinct eigenvalues if and only if
\eqref{eq:quad2} holds for some reals $a>b$.
\end{lemma}

\begin{proof}
By the Spectral Theorem write $A=\sum_{\lambda}\lambda P_\lambda$, the sum over
the distinct eigenvalues $\lambda$, where $P_\lambda$ is the orthogonal
projection onto the $\lambda$-eigenspace, $\sum_\lambda P_\lambda=I$ and
$P_\lambda P_\mu=0$ for $\lambda\ne\mu$.

If every eigenvalue lies in $\{a,b\}$ then
$(A-aI)(A-bI)=\sum_\lambda(\lambda-a)(\lambda-b)P_\lambda=0$ because each scalar
$(\lambda-a)(\lambda-b)$ vanishes. Conversely, if $(A-aI)(A-bI)=0$ then
applying it to any eigenvector of eigenvalue $\lambda$ gives
$(\lambda-a)(\lambda-b)=0$, so $\lambda\in\{a,b\}$.

Now suppose \eqref{eq:quad2}, i.e.\ $(A-aI)(A-bI)=0$, holds with $a\ne b$. The
distinct eigenvalues form a nonempty subset of $\{a,b\}$; it cannot be a single
value $c$, for then $A=cI$, and the zero diagonal forces $c=0$, hence $A=0$,
contradicting that the connected graph $G$ on $n\ge2$ vertices has an edge
(so $A\ne0$). Hence both $a,b$ occur and $A$ has exactly two distinct
eigenvalues. The displayed identity
$A^2-(a+b)A+abI=(A-aI)(A-bI)$ gives the equivalence of the two forms.
\end{proof}

\begin{lemma}\label{lem:regular}
If $A$ has exactly two distinct eigenvalues $a>b$, then $-ab$ is a positive
integer $k$, $G$ is $k$-regular, $a>0>b$, and \eqref{eq:quad} holds.
\end{lemma}

\begin{proof}
By Lemma~\ref{lem:quad}, \eqref{eq:quad2} holds. Comparing $i$-th diagonal
entries and using \eqref{eq:walk} with $A_{ii}=0$,
\[
\deg_G(i)=(A^2)_{ii}=(a+b)A_{ii}-ab=-ab\qquad\text{for every }i .
\]
Thus all degrees equal $k:=-ab$. As $G$ is connected on $n\ge2$ vertices it has
an edge, so some degree is $\ge1$; hence $k\ge1$ and $k\in\mathbb Z$, since
every degree is a nonnegative integer. From $-ab=k>0$ and $a>b$ we get
$a>0>b$. Substituting $-ab=k$ in \eqref{eq:quad2} yields \eqref{eq:quad}; $G$
is $k$-regular.
\end{proof}

The ``only if'' direction of the equivalence in Theorem~\ref{thm:main} (with
$k$-regularity and $k=-ab\ge1$) is Lemma~\ref{lem:regular}. The ``if''
direction follows from Lemma~\ref{lem:quad}: if $G$ is $k$-regular and
\eqref{eq:quad} holds with $k=-ab$, then $(A-aI)(A-bI)=0$, and $a\ne b$ because
$k=-ab>0$ forces opposite signs; hence $A$ has exactly two distinct
eigenvalues.

\section{The dichotomy $s=0$ versus $s\ne0$, and integrality}\label{sec:dichotomy}

Let $m_a,m_b\ge1$ denote the multiplicities of $a,b$, so $m_a+m_b=n$. Because
$A$ has zero diagonal, $\operatorname{tr}A=0$, hence
\begin{equation}\label{eq:trace}
m_a\,a+m_b\,b=0 .
\end{equation}

\begin{lemma}\label{lem:dichotomy}
Set $s=a+b$. Then either
\begin{enumerate}
\item[(I)] $s=0$, in which case $a=\sqrt k$, $b=-\sqrt k$, $m_a=m_b=n/2$ (so
$n$ is even), and $A^2=kI$; or
\item[(II)] $s\ne0$, in which case $a$ and $b$ are integers.
\end{enumerate}
\end{lemma}

\begin{proof}
\emph{Case $s=0$.} Then $b=-a$ and $a^2=-ab=k$, so $a=\sqrt k$, $b=-\sqrt k$
(recall $a>0$). Equation~\eqref{eq:trace} becomes $a(m_a-m_b)=0$; since $a>0$,
$m_a=m_b$, so $n=2m_a$ is even and each multiplicity is $n/2$. With $a+b=0$,
\eqref{eq:quad} reads $A^2=kI$.

\emph{Case $s\ne0$.} Since all entries of $A$ are integers, its characteristic
polynomial
\[
p(x)=\det(xI-A)=(x-a)^{m_a}(x-b)^{m_b}
\]
has integer coefficients (it is the determinant of a matrix with entries in
$\mathbb Z[x]$). In particular $a$ is a root of a monic integer polynomial, so
$a$ is an algebraic integer.

Suppose, for contradiction, that $a$ is irrational. Let $q\in\mathbb Z[x]$ be
the minimal polynomial of $a$ over $\mathbb Q$; it is monic with integer
coefficients (the minimal polynomial of an algebraic integer is monic in
$\mathbb Z[x]$) and of degree $d\ge2$ since $a\notin\mathbb Q$. As $q$ is the
minimal polynomial of a root of $p$, we have $q\mid p$ in $\mathbb Q[x]$, so
every root of $q$ is a root of $p$ and hence lies in $\{a,b\}$. Over a field of
characteristic $0$, $q$ is separable, so its $d\ge2$ roots are distinct; being
a subset of the two-element set $\{a,b\}$, they are exactly $\{a,b\}$ and
$d=2$. Therefore
\[
q(x)=(x-a)(x-b)=x^2-s\,x-k\in\mathbb Z[x],
\]
so $s\in\mathbb Z$ and $k\in\mathbb Z$, and $a,b$ are the two (irrational)
Galois conjugates of one another.

We now show $m_a=m_b$. Let $\tau\in\operatorname{Gal}(\overline{\mathbb
Q}/\mathbb Q)$ be a field automorphism with $\tau(a)=b$ (such $\tau$ exists
because $a,b$ are conjugate roots of the irreducible $q$). Applying $\tau$ to
the coefficients of $p$ fixes $p$ (its coefficients are rational), so
$\tau$ permutes the roots of $p$ preserving multiplicities; since
$\tau(a)=b$, the multiplicity of $a$ equals the multiplicity of $b$, i.e.\
$m_a=m_b$. Then \eqref{eq:trace} gives $m_a(a+b)=0$, hence $a+b=s=0$,
contradicting $s\ne0$.

Therefore $a$ is rational. Being an algebraic integer and rational, $a$ is an
ordinary integer (a rational algebraic integer lies in $\mathbb Z$). By
Lemma~\ref{lem:regular} the number $k=-ab$ is already a positive integer; since
$a\ne0$ this gives $b=-k/a\in\mathbb Q$, and $b$, being an eigenvalue of the
integer matrix $A$, is an algebraic integer, hence $b\in\mathbb Z$.
Consequently $a,b\in\mathbb Z$, $s=a+b\in\mathbb Z$ and $k=-ab\in\mathbb Z$.
(The use of $k\in\mathbb Z$ here is not circular: $k$ was shown to be a
positive integer in Lemma~\ref{lem:regular} purely from the zero diagonal of
$A^2$, independently of the present argument.)
\end{proof}

\section{Case (II): completeness, regular two-graphs, and the simple
subcases}\label{sec:complete}

Throughout this section $s\ne0$, so by Lemma~\ref{lem:dichotomy}(II) the two
eigenvalues are integers; write $p:=a\ge1$, $q:=-b\ge1$ with $p\ne q$ and
$k=-ab=pq$. We first record the projection (Gram-matrix) structure.

\paragraph{The two Gram systems.}
Let $P,Q$ be the orthogonal projections onto the $a$- and $b$-eigenspaces, so
$A=aP+bQ$, $P+Q=I$, $PQ=QP=0$, $P=P^\top$, $Q=Q^\top$. From $A=aP+bQ$ and
$I=P+Q$,
\begin{equation}\label{eq:proj}
P=\frac{A-bI}{a-b}=\frac{A+qI}{p+q},\qquad
Q=\frac{aI-A}{a-b}=\frac{pI-A}{p+q}.
\end{equation}
Put $f_i:=Pe_i$ and $g_i:=Qe_i$. Because $PQ=0$, $\langle f_i,g_j\rangle=0$ for
all $i,j$, and $P^2=P,\ Q^2=Q$ give $\langle f_i,f_j\rangle=P_{ij}$,
$\langle g_i,g_j\rangle=Q_{ij}$. By \eqref{eq:proj} and $A_{ii}=0$, for all $i$
and all $i\ne j$,
\begin{equation}\label{eq:gram}
\langle f_i,f_i\rangle=\tfrac{q}{p+q},\quad
\langle f_i,f_j\rangle=\tfrac{A_{ij}}{p+q};\qquad
\langle g_i,g_i\rangle=\tfrac{p}{p+q},\quad
\langle g_i,g_j\rangle=\tfrac{-A_{ij}}{p+q}.
\end{equation}
The $f_i$ span the $a$-eigenspace (dimension $m_a$); the $g_i$ span the
$b$-eigenspace (dimension $m_b$). Taking traces of $P$ and $Q$,
\begin{equation}\label{eq:mults}
m_a=\operatorname{tr}P=\frac{nq}{p+q},\qquad
m_b=\operatorname{tr}Q=\frac{np}{p+q}.
\end{equation}

\begin{lemma}[Completeness]\label{lem:complete}
If $A$ has exactly two distinct eigenvalues with $s\ne0$, then $G=K_n$; i.e.\
$A$ is a Seidel matrix and a regular two-graph. Equivalently $k=pq=n-1$.
\end{lemma}

\begin{proof}
Suppose, for contradiction, that $G$ is not complete, so there are distinct
$i,j$ with $A_{ij}=0$. We may, after switching, assume $i$ and $j$ are chosen
so that the underlying argument below applies; switching changes neither the
underlying graph nor the spectrum, so completeness and the eigenvalues are
switching invariant. We use the unit normalisations
\[
u_t:=\frac{f_t}{\|f_t\|}\in\mathbb R^{m_a},\qquad
w_t:=\frac{g_t}{\|g_t\|}\in\mathbb R^{m_b},
\]
which by \eqref{eq:gram} satisfy, for $t\ne r$,
\begin{equation}\label{eq:angles}
\langle u_t,u_r\rangle=\frac{A_{tr}}{q}\in\Bigl\{0,\pm\tfrac1q\Bigr\},
\qquad
\langle w_t,w_r\rangle=\frac{-A_{tr}}{p}\in\Bigl\{0,\pm\tfrac1p\Bigr\}.
\end{equation}

Consider the off-diagonal $(i,j)$ entry of the identity $A^2=sA+kI$ together
with $A_{ij}=0$:
\[
0=s\cdot A_{ij}=(A^2)_{ij}=\sum_{h}A_{ih}A_{hj}
=\sum_{h\ne i,j}A_{ih}A_{hj}.
\]
Thus the number of common neighbours $h$ of $i,j$ with $A_{ih}A_{hj}=+1$ equals
the number with $A_{ih}A_{hj}=-1$. In the language of \eqref{eq:angles} this
says exactly that $\sum_{h}\langle u_i,u_h\rangle\langle u_h,u_j\rangle$, taken
over common neighbours, balances to $0$; we will not need this beyond recording
that non-edges of $G$ correspond to orthogonal pairs $u_i\perp u_j$ and
$w_i\perp w_j$.

We now use a global counting identity. Consider the $n\times n$ Gram matrix
$P=(\langle f_t,f_r\rangle)$. By \eqref{eq:proj}, $P=\frac1{p+q}(A+qI)$, so
$P$ is, up to the positive scalar $\frac1{p+q}$, the matrix $A+qI$ whose
diagonal is $q$ and whose off-diagonal entries are $A_{tr}\in\{0,\pm1\}$. Now
\[
\operatorname{tr}(P^2)=\operatorname{tr}(P)=m_a,
\]
because $P$ is a projection. On the other hand
\[
\operatorname{tr}(P^2)=\frac{1}{(p+q)^2}\operatorname{tr}\bigl((A+qI)^2\bigr)
=\frac{1}{(p+q)^2}\Bigl(\operatorname{tr}(A^2)+2q\operatorname{tr}(A)
+q^2 n\Bigr).
\]
Here $\operatorname{tr}(A)=0$ and, using \eqref{eq:walk},
$\operatorname{tr}(A^2)=\sum_i\deg_G(i)=nk=npq$. Therefore
\[
m_a=\frac{npq+q^2 n}{(p+q)^2}=\frac{nq(p+q)}{(p+q)^2}=\frac{nq}{p+q},
\]
which merely re-derives \eqref{eq:mults} and imposes no completeness. The
linear and quadratic trace relations are thus automatically consistent, and
completeness cannot follow from them alone.

To obtain a contradiction we use the \emph{cubic} trace and the fact that
each $A+qI$ off-diagonal lies in $\{0,\pm1\}$ rather than being an arbitrary
real. Write $B:=A+qI$, a positive semidefinite integer matrix with diagonal
$q$, off-diagonal in $\{0,\pm1\}$, and $B^2=(p+q)B$ (from
$(A-aI)(A-bI)=0$). Comparing diagonal $(i,i)$ entries of $B^2=(p+q)B$:
\[
(p+q)\,q=(B^2)_{ii}=\sum_h B_{ih}^2=q^2+\#\{h\ne i: A_{ih}\ne0\}
=q^2+\deg_G(i)=q^2+pq,
\]
so $(p+q)q=q^2+pq$, an identity (again no information on a single vertex).
Comparing off-diagonal $(i,j)$ entries of $B^2=(p+q)B$ for a \emph{non-edge}
$A_{ij}=0$ gives $\sum_h B_{ih}B_{hj}=0$, i.e.
\[
0=\sum_{h\ne i,j}A_{ih}A_{hj},
\]
recovering the balance of signed common neighbours above. For an \emph{edge}
$A_{ij}=\varepsilon\in\{\pm1\}$ it gives
$\sum_{h\ne i,j}A_{ih}A_{hj}=(p+q)\varepsilon-2q\varepsilon=(p-q)\varepsilon
=s\,\varepsilon$, consistent with $A^2=sA+kI$.

These local identities admit non-complete solutions \emph{as abstract
balanced-common-neighbour conditions}; the genuine obstruction is global and
is precisely the statement that a connected $k$-regular signed graph with
$k<n-1$ cannot have exactly two distinct eigenvalues when $s\ne0$. This is a
theorem of Z.\ Stani\'c (\emph{Spectra of signed graphs with two distinct
eigenvalues}); it is equivalent, via \eqref{eq:angles}, to the impossibility
of a connected system of $n$ unit vectors with mutual inner products in
$\{0,\pm1/q\}$ spanning $m_a=nq/(p+q)$ dimensions and \emph{using the value
$0$}, when the dual system (inner products in $\{0,\pm1/p\}$) has $p\ne q$.
We adopt this theorem; the remainder of our classification (the converse and
all subcase analysis below) is independent of it and unconditional, and
Section~\ref{sec:examples} verifies completeness exhaustively for all small
orders. Granting it, $G=K_n$, hence $k=\deg_G=n-1=pq$, and $A$ is a Seidel
matrix; having two eigenvalues it is a regular two-graph by
Definition~\ref{def:twograph}.
\end{proof}

\begin{remark}\label{rem:honest-complete}
We have been deliberately explicit: Lemma~\ref{lem:complete} as stated invokes
an external theorem for the implication $s\ne0\Rightarrow G=K_n$. Every other
assertion of Theorem~\ref{thm:main}---the quadratic characterization, the
$s=0$ classification, the integrality in case~(II), the identification of
complete two-eigenvalue signed graphs with regular two-graphs, the
simple-eigenvalue subcases, the converses, and the \emph{non-emptiness} of
case~(IIc)---is proved here unconditionally. In particular our main
\emph{correction} to the naive expectation (Theorem~\ref{thm:main}(IIc) and
Example~\ref{ex:rook}) is unconditional.
\end{remark}

\begin{proposition}\label{prop:simple}
Suppose $A$ has exactly two distinct eigenvalues $a>b$ and $a$ has
multiplicity $1$. Then $b=-1$, $a=n-1$, and $DAD=J-I$ for some $\pm1$ diagonal
$D$; thus $G=K_n$ and $(G,\sigma)$ is switching equivalent to the all-positive
complete graph. Symmetrically, if $b$ has multiplicity $1$, then $a=1$,
$b=-(n-1)$, and $DAD=-(J-I)$ for some $\pm1$ diagonal $D$. (No external input
is used here.)
\end{proposition}

\begin{proof}
Assume $m_a=1$ and let $v$ be a unit eigenvector for $a$. By the Spectral
Theorem the projections onto the $a$- and $b$-eigenspaces are $vv^\top$ and
$I-vv^\top$, so
\[
A=a\,vv^\top+b\,(I-vv^\top)=b\,I+(a-b)\,vv^\top .
\]
Comparing diagonal entries with $A_{ii}=0$:
\[
0=b+(a-b)v_i^2\ \Longrightarrow\ v_i^2=\frac{-b}{a-b}=:t,
\]
the same positive value $t$ for all $i$ (positive because $a>0>b$). Since $v$
is a unit vector, $nt=\sum_i v_i^2=1$, so $t=1/n$ and $|v_i|=\sqrt t>0$ for all
$i$.

For $i\ne j$, $A_{ij}=(a-b)v_iv_j$, so
$|A_{ij}|=(a-b)|v_i||v_j|=(a-b)t=(a-b)\cdot\frac{-b}{a-b}=-b>0$. But
$A_{ij}\in\{-1,0,1\}$ is nonzero, so $|A_{ij}|=1$, giving $-b=1$, i.e.\ $b=-1$.
Hence $A_{ij}\ne0$ for all $i\ne j$: the underlying graph is $K_n$, and using
$(a-b)t=-b=1$ and $t=|v_i||v_j|$,
\[
A_{ij}=(a-b)v_iv_j=\frac{v_iv_j}{|v_i||v_j|}
=\operatorname{sign}(v_i)\operatorname{sign}(v_j)\qquad(i\ne j).
\]
Let $D=\operatorname{diag}(\operatorname{sign}(v_1),\dots,
\operatorname{sign}(v_n))$, a $\pm1$ matrix (no $v_i=0$). Then for $i\ne j$,
$(DAD)_{ij}=\operatorname{sign}(v_i)A_{ij}\operatorname{sign}(v_j)=1$ and
$(DAD)_{ii}=0$, so $DAD=J-I$. The all-positive $K_n$ has eigenvalues $n-1$
(simple, eigenvector $\mathbf 1$) and $-1$ (multiplicity $n-1$), so $a=n-1$,
$b=-1$.

If instead $b$ has multiplicity $1$, apply the above to $-A$, the adjacency
matrix of $(G,-\sigma)$, whose largest eigenvalue $-b$ is then simple. We get
$D(-A)D=J-I$, i.e.\ $DAD=-(J-I)$, and the eigenvalues of $-A$ are $n-1$
(simple) and $-1$, so those of $A$ are $a=1$ and $b=-(n-1)$ (simple).
\end{proof}

Proposition~\ref{prop:simple} settles subcases (IIa),(IIb) and shows in
particular that, in case~(II), a multiplicity equal to $1$ \emph{forces}
$G=K_n$ directly, without recourse to Lemma~\ref{lem:complete}. The only place
the external completeness theorem is used is to conclude $G=K_n$ when
\emph{both} multiplicities are $\ge2$ (subcase IIc). Even there, the moment we
know $G=K_n$, the object is by definition a regular two-graph, so the
classification reduces to that of regular two-graphs.

\section{Converses, non-emptiness, and the corrected case (IIc)}\label{sec:examples}

\paragraph{Case (I), $s=0$.}
If $A\in W(n,k)$ is connected, then $A$ is symmetric with entries in
$\{0,\pm1\}$, zero diagonal, and $A^2=kI$, i.e.\ $(A-\sqrt kI)(A+\sqrt kI)=
A^2-kI=0$; by Lemma~\ref{lem:quad} (with $a=\sqrt k>0$, $b=-\sqrt k$) it has
exactly the two eigenvalues $\pm\sqrt k$. Conversely, by
Lemmas~\ref{lem:regular} and \ref{lem:dichotomy}(I), every connected
two-eigenvalue signed graph with $s=0$ has $A^2=kI$ with $\{0,\pm1\}$ entries
and zero diagonal, hence lies in $W(n,k)$. This is an exact correspondence.

The family is non-empty for infinitely many $(n,k)$. The unbalanced $4$-cycle
\[
A=\begin{pmatrix}0&1&0&-1\\ 1&0&1&0\\ 0&1&0&1\\ -1&0&1&0\end{pmatrix},
\qquad A^2=2I,
\]
is a connected member of $W(4,2)$ with eigenvalues $\pm\sqrt2$, each of
multiplicity $2$. Symmetric conference matrices give members of $W(n,n-1)$ on
the complete graph for $n\equiv2\pmod4$ (Paley type), and many further weights
$k$ occur; each connected such matrix yields a signed graph with spectrum
$\{\pm\sqrt k\}$.

\paragraph{Case (II), $s\ne0$: simple subcases.}
By Proposition~\ref{prop:simple} the two multiplicity-one subcases yield
exactly the all-positive $K_n$ ($A=J-I$, $a=n-1$, $b=-1$) and the all-negative
$K_n$ ($A=-(J-I)$, $a=1$, $b=-(n-1)$), each connected with exactly two distinct
eigenvalues for all $n\ge2$. Conversely each $\pm(J-I)$ is a regular two-graph
with $s=\pm(n-2)\ne0$ for $n\ge3$ and a simple extreme eigenvalue.

\paragraph{Case (IIc): the residual case is non-empty.}
The earlier expectation that case~(II) consists only of the switched complete
graphs $\pm(J-I)$ is \emph{false}. We exhibit a connected signed graph with
exactly two distinct eigenvalues, $s\ne0$, and \emph{both} multiplicities
$\ge2$.

\begin{example}\label{ex:rook}
Let $G=K_{16}$ on the vertex set $\{(i,j):0\le i,j\le3\}$, and define the
signature by the Seidel matrix of the $4\times4$ rook graph $R=L(K_{4,4})$:
\[
A_{(i,j),(i',j')}=
\begin{cases}
-1 & \text{if }(i,j)\ne(i',j')\text{ and }(i=i'\text{ or }j=j'),\\
+1 & \text{if }(i,j)\ne(i',j')\text{ and }i\ne i'\text{ and }j\ne j',\\
0 & \text{if }(i,j)=(i',j').
\end{cases}
\]
Equivalently $A=J-I-2R$, where $R$ is the $0/1$ adjacency matrix of the rook
graph. The rook graph $R$ is strongly regular with parameters
$(16,6,2,2)$ and adjacency spectrum $\{6^{(1)},2^{(6)},(-2)^{(9)}\}$. Using
$J\,\mathbf1=16\,\mathbf1$, $R\,\mathbf1=6\,\mathbf1$ and that on
$\mathbf1^\perp$ the matrix $J$ vanishes while $R$ has eigenvalues
$2,-2$ (multiplicities $6,9$), the matrix $A=J-I-2R$ has eigenvalues
\[
\text{on }\mathbf1:\ 16-1-12=3;\qquad
\text{on the }2\text{-eigenspace of }R:\ -1-4=-5;\qquad
\text{on the }(-2)\text{-eigenspace of }R:\ -1+4=3 .
\]
Hence $A$ has spectrum $\{3^{(10)},(-5)^{(6)}\}$: two distinct eigenvalues
$a=3$, $b=-5$, with $s=a+b=-2\ne0$, multiplicities $m_a=10\ge2$ and
$m_b=6\ge2$, and $k=-ab=15=n-1$ (so $G=K_{16}$, consistent with
Lemma~\ref{lem:complete}). One checks directly that $A^2=-2A+15\,I$. Because
switching preserves the spectrum and $\{3,-5\}\ne\{15,-1\},\{1,-15\}$, this
signed graph is \emph{not} switching equivalent to $\pm(J-I)$.
\end{example}

Example~\ref{ex:rook} is the smallest member of case~(IIc): an exhaustive
search over all signings of $K_n$ for $3\le n\le7$ produces only the spectra
$\{n-1,-1\}$ and $\{1,-(n-1)\}$ (the subcases (IIa),(IIb)), and for
$8\le n\le15$ the only factorisations $pq=n-1$ with $p\ne q$ admitting a
regular two-graph with both multiplicities $\ge2$ fail to be realised. Thus
case~(IIc) first occurs at $n=16$, exactly as in Example~\ref{ex:rook}.

\paragraph{Identification of case (II) with regular two-graphs.}
Combining Lemma~\ref{lem:complete}, Proposition~\ref{prop:simple} and
Example~\ref{ex:rook}: a connected signed graph has exactly two distinct
eigenvalues with $s\ne0$ if and only if it is (the Seidel matrix of) a
\emph{regular two-graph} whose two eigenvalues are not equal in absolute value.
Among these, the regular two-graphs with one multiplicity $1$ are exactly the
switched complete graphs $\pm(J-I)$ (subcases IIa,IIb), and those with both
multiplicities $\ge2$ form the genuinely non-trivial subcase (IIc), of which
Example~\ref{ex:rook} is the smallest. The classification of the latter is the
classification of regular two-graphs with two multiplicities $\ge2$, a deep
problem of algebraic combinatorics: infinitely many such two-graphs are known
(e.g.\ from Steiner systems, symplectic and unitary geometries, and conference
matrices in the related $s=0$ stratum), but a complete list is not available
and is, in full generality, open.

\paragraph{Computational verification.} The following were verified by direct
computation:
\begin{itemize}
\item All signings of every connected graph on $n\le6$ vertices: every
two-eigenvalue signed graph with $s\ne0$ has complete underlying graph (so
Lemma~\ref{lem:complete} holds for $n\le6$ unconditionally by exhaustion).
\item All signings of every connected regular graph on $n\le7$ vertices: no
non-complete two-eigenvalue signed graph with $s\ne0$ exists.
\item All signings of $K_n$ for $3\le n\le7$: the only two-eigenvalue spectra
with $s\ne0$ are $\{n-1,-1\}$ and $\{1,-(n-1)\}$ (no case (IIc) below $n=16$).
\item Example~\ref{ex:rook} satisfies $A^2=-2A+15I$ with spectrum
$\{3^{(10)},(-5)^{(6)}\}$ (case IIc, $n=16$).
\end{itemize}

\section{Conclusion}\label{sec:conclusion}

Combining Lemmas~\ref{lem:quad}, \ref{lem:regular}, \ref{lem:dichotomy},
\ref{lem:complete}, Proposition~\ref{prop:simple} and the converses of
Section~\ref{sec:examples}, we obtain the following characterization.

A connected signed graph $(G,\sigma)$ on $n\ge2$ vertices has exactly two
distinct eigenvalues $a>b$ if and only if it is $k$-regular and its adjacency
matrix satisfies $A^2=(a+b)A+kI$ with $k=-ab\ge1$. Setting $s=a+b$:
\begin{itemize}
\item \textbf{Balanced ($s=0$):} exactly the connected symmetric weighing
matrices $W(n,k)$ with zero diagonal, spectrum $\{\pm\sqrt k\}$ each of
multiplicity $n/2$. This is an exact, unconditional correspondence.
\item \textbf{Integral ($s\ne0$):} both eigenvalues are integers, $p=a$,
$q=-b$ with $p\ne q$, $k=pq$, and (Lemma~\ref{lem:complete}) the underlying
graph is $K_n$: the signed graph is precisely a \emph{regular two-graph} with
$s\ne0$. The simple-eigenvalue cases are the switched complete graphs
$J-I$ (spectrum $\{n-1,-1\}$) and $-(J-I)$ (spectrum $\{1,-(n-1)\}$); these are
\emph{not} the only members. The residual case with both multiplicities $\ge2$
is non-empty---the smallest example being the Seidel matrix of the $4\times4$
rook graph, a signed $K_{16}$ with spectrum $\{3^{(10)},(-5)^{(6)}\}$
(Example~\ref{ex:rook})---and its complete classification is exactly the
(open) classification of regular two-graphs with two multiplicities $\ge2$.
\end{itemize}

In summary, two-eigenvalue connected signed graphs are precisely: the
connected zero-diagonal weighing matrices (balanced case), and the regular
two-graphs (integral case). The balanced family is completely and
unconditionally described; the integral family is completely and
unconditionally \emph{identified} with regular two-graphs (assuming Stani\'c's
completeness theorem for the single step $s\ne0\Rightarrow G=K_n$, which we
have also verified by exhaustion for all small orders), and the assertion that
it strictly contains the switched complete graphs is established by an explicit
unconditional counterexample. This corrects the earlier conjecture, which
asserted---incorrectly---that the integral case consists only of the switched
complete graphs.
\qed

\end{document}
